% --- Archon physics formalization source begin ---
% archon:physics
% archon:covers PhyXMiniProblems/problem_phyx_mini_0172.lean
% archon:source-report reports/phyx_mini/problem_phyx_mini_0172.source.json

\chapter{Physics problem phyx\_mini\_0172}
\label{ch:PhyXMiniProblems_problem_phyx_mini_0172}

\paragraph{Problem source.}
\#\# Physical scenario

A uniform \textbackslash{}( 165\textbackslash{},\textbackslash{}text\{N\} \textbackslash{}) bar is supported horizontally by two identical wires \textbackslash{}( A \textbackslash{}) and \textbackslash{}( B \textbackslash{}) (\textbackslash{}textbf\{figure\}). A small \textbackslash{}( 185\textbackslash{},\textbackslash{}text\{N\} \textbackslash{}) cube of lead is placed three-fourths of the way from \textbackslash{}( A \textbackslash{}) to \textbackslash{}( B \textbackslash{}). The wires are each \textbackslash{}( 75.0\textbackslash{},\textbackslash{}text\{cm\} \textbackslash{}) long and have a mass of \textbackslash{}( 5.50\textbackslash{},\textbackslash{}text\{g\} \textbackslash{}).

\#\# Question

If both of them are simultaneously plucked at the center, what is the frequency of the beats that they will produce when vibrating in their fundamental?

\#\# Answer choices

- A: A: 27.3Hz

- B: B: 38.0Hz

- C: C: 89.6Hz

- D: D: 38.5Hz

\#\# Figure caption (auxiliary; use the image as primary evidence)

The image shows a simple illustration of a room-like structure with three vertical black line segments outlining the sides. The top line segment has a shaded upper edge resembling a ceiling. Inside the room, there is a horizontal rectangular shape labeled "Bar," extending across the bottom, depicted with a brown color and shaded to suggest three dimensions.

Near the right side of the bar, there is a small brown cube labeled "Cube." The letters "A" and "B" are positioned vertically on the left and right sides of the room, suggesting measurements or positions along the walls. The overall layout conveys a basic mechanical or architectural diagram with labeled components.

\#\# Dataset metadata

Resonance and Harmonics; Physical Model Grounding Reasoning, Multi-Formula Reasoning

\paragraph{Recorded answer/context.}
A: A: 27.3Hz

\paragraph{Figure/image path.}
/root/proposal\_for\_physic/science-mango/phyx\_mini\_run/phyx\_data/test\_image/172.png

\paragraph{Formalization target.}
create a compiling Lean file with sorry bodies at `PhyXMiniProblems/problem\_phyx\_mini\_0172.lean`. The Lean declarations must preserve the physical quantities, dimensions or dimensional roles, figure labels, governing-law hypotheses, and final relation expressed by this problem.
Use Mathlib/Physlib names found through LeanExplore where available. If a physics API is missing, introduce faithful local abstractions rather than scalar placeholder aliases.

\begin{theorem}[Physics formalization target]
\label{thm:physics:phyx_mini_0172:target}
\lean{PhyXMiniProblems.ProblemPhyXMini0172.problem_phyx_mini_0172}
\uses{def:physics:phyx-mini-0172:phyxminiproblems-problemphyxmini0172-frequencyinhertz, def:physics:phyx-mini-0172:phyxminiproblems-problemphyxmini0172-suspendedbarwiresetup, def:physics:phyx-mini-0172:phyxminiproblems-problemphyxmini0172-matchesprimaryfigure, def:physics:phyx-mini-0172:phyxminiproblems-problemphyxmini0172-matchesbarandsupportgeometry, def:physics:phyx-mini-0172:phyxminiproblems-problemphyxmini0172-matchesstatedmeasurements, def:physics:phyx-mini-0172:phyxminiproblems-problemphyxmini0172-matchessimultaneousfundamentalpluck, def:physics:phyx-mini-0172:phyxminiproblems-problemphyxmini0172-hasphysicalparameters, def:physics:phyx-mini-0172:phyxminiproblems-problemphyxmini0172-satisfiesstaticequilibrium, def:physics:phyx-mini-0172:phyxminiproblems-problemphyxmini0172-stretchedstringfundamentalhertz, def:physics:phyx-mini-0172:phyxminiproblems-problemphyxmini0172-satisfiesstretchedstringfundamentallaw, def:physics:phyx-mini-0172:phyxminiproblems-problemphyxmini0172-satisfiesbeatfrequencylaw, def:physics:phyx-mini-0172:phyxminiproblems-problemphyxmini0172-isclosestanswerchoice}
The assigned autoformalize agent should translate this physics problem into Lean declarations in the covered file, with theorem and lemma proof bodies written as `by sorry`.
\end{theorem}
\begin{proof}
This is an autoformalization task, not a proof task. Produce statements that can later be proved by the physics prover without weakening the physical model.
\end{proof}

% --- Archon declaration topology begin ---
\section{Declaration topology}
\begin{definition}[force Dimension]
  \label{def:physics:phyx-mini-0172:phyxminiproblems-problemphyxmini0172-forcedimension}
  \lean{PhyXMiniProblems.ProblemPhyXMini0172.forceDimension}
  The physical dimension of force, mass * length / time\textasciicircum{}2
\end{definition}
\begin{proof}
  This is the declared model definition of force Dimension.
\end{proof}
\begin{definition}[Length Quantity]
  \label{def:physics:phyx-mini-0172:phyxminiproblems-problemphyxmini0172-lengthquantity}
  \lean{PhyXMiniProblems.ProblemPhyXMini0172.LengthQuantity}
  A unit-independent physical length
\end{definition}
\begin{proof}
  This is the declared model definition of Length Quantity.
\end{proof}
\begin{definition}[Mass Quantity]
  \label{def:physics:phyx-mini-0172:phyxminiproblems-problemphyxmini0172-massquantity}
  \lean{PhyXMiniProblems.ProblemPhyXMini0172.MassQuantity}
  A unit-independent physical mass
\end{definition}
\begin{proof}
  This is the declared model definition of Mass Quantity.
\end{proof}
\begin{definition}[Force Quantity]
  \label{def:physics:phyx-mini-0172:phyxminiproblems-problemphyxmini0172-forcequantity}
  \lean{PhyXMiniProblems.ProblemPhyXMini0172.ForceQuantity}
  \uses{def:physics:phyx-mini-0172:phyxminiproblems-problemphyxmini0172-forcedimension}
  A unit-independent physical force, used for weights and wire tensions
\end{definition}
\begin{proof}
  This is the declared model definition of Force Quantity.
\end{proof}
\begin{definition}[Time Quantity]
  \label{def:physics:phyx-mini-0172:phyxminiproblems-problemphyxmini0172-timequantity}
  \lean{PhyXMiniProblems.ProblemPhyXMini0172.TimeQuantity}
  A unit-independent physical time
\end{definition}
\begin{proof}
  This is the declared model definition of Time Quantity.
\end{proof}
\begin{definition}[Frequency Quantity]
  \label{def:physics:phyx-mini-0172:phyxminiproblems-problemphyxmini0172-frequencyquantity}
  \lean{PhyXMiniProblems.ProblemPhyXMini0172.FrequencyQuantity}
  A unit-independent physical frequency, carrying inverse-time dimension
\end{definition}
\begin{proof}
  This is the declared model definition of Frequency Quantity.
\end{proof}
\begin{definition}[quantity Readout]
  \label{def:physics:phyx-mini-0172:phyxminiproblems-problemphyxmini0172-quantityreadout}
  \lean{PhyXMiniProblems.ProblemPhyXMini0172.quantityReadout}
  Read a dimension-tagged physical quantity in a selected system of units
\end{definition}
\begin{proof}
  This is the declared model definition of quantity Readout.
\end{proof}
\begin{definition}[length In Meters]
  \label{def:physics:phyx-mini-0172:phyxminiproblems-problemphyxmini0172-lengthinmeters}
  \lean{PhyXMiniProblems.ProblemPhyXMini0172.lengthInMeters}
  \uses{def:physics:phyx-mini-0172:phyxminiproblems-problemphyxmini0172-lengthquantity, def:physics:phyx-mini-0172:phyxminiproblems-problemphyxmini0172-quantityreadout}
  Length readout in metres
\end{definition}
\begin{proof}
  This is the declared model definition of length In Meters.
\end{proof}
\begin{definition}[length In Centimeters]
  \label{def:physics:phyx-mini-0172:phyxminiproblems-problemphyxmini0172-lengthincentimeters}
  \lean{PhyXMiniProblems.ProblemPhyXMini0172.lengthInCentimeters}
  \uses{def:physics:phyx-mini-0172:phyxminiproblems-problemphyxmini0172-lengthquantity, def:physics:phyx-mini-0172:phyxminiproblems-problemphyxmini0172-quantityreadout}
  Length readout in centimetres, the unit used for each wire in the source
\end{definition}
\begin{proof}
  This is the declared model definition of length In Centimeters.
\end{proof}
\begin{definition}[mass In Kilograms]
  \label{def:physics:phyx-mini-0172:phyxminiproblems-problemphyxmini0172-massinkilograms}
  \lean{PhyXMiniProblems.ProblemPhyXMini0172.massInKilograms}
  \uses{def:physics:phyx-mini-0172:phyxminiproblems-problemphyxmini0172-massquantity, def:physics:phyx-mini-0172:phyxminiproblems-problemphyxmini0172-quantityreadout}
  Mass readout in kilograms
\end{definition}
\begin{proof}
  This is the declared model definition of mass In Kilograms.
\end{proof}
\begin{definition}[mass In Grams]
  \label{def:physics:phyx-mini-0172:phyxminiproblems-problemphyxmini0172-massingrams}
  \lean{PhyXMiniProblems.ProblemPhyXMini0172.massInGrams}
  \uses{def:physics:phyx-mini-0172:phyxminiproblems-problemphyxmini0172-massquantity, def:physics:phyx-mini-0172:phyxminiproblems-problemphyxmini0172-quantityreadout}
  Mass readout in grams, the unit used for each wire in the source
\end{definition}
\begin{proof}
  This is the declared model definition of mass In Grams.
\end{proof}
\begin{definition}[force In Newtons]
  \label{def:physics:phyx-mini-0172:phyxminiproblems-problemphyxmini0172-forceinnewtons}
  \lean{PhyXMiniProblems.ProblemPhyXMini0172.forceInNewtons}
  \uses{def:physics:phyx-mini-0172:phyxminiproblems-problemphyxmini0172-forcequantity, def:physics:phyx-mini-0172:phyxminiproblems-problemphyxmini0172-quantityreadout}
  Force readout in SI newtons
\end{definition}
\begin{proof}
  This is the declared model definition of force In Newtons.
\end{proof}
\begin{definition}[time In Seconds]
  \label{def:physics:phyx-mini-0172:phyxminiproblems-problemphyxmini0172-timeinseconds}
  \lean{PhyXMiniProblems.ProblemPhyXMini0172.timeInSeconds}
  \uses{def:physics:phyx-mini-0172:phyxminiproblems-problemphyxmini0172-timequantity, def:physics:phyx-mini-0172:phyxminiproblems-problemphyxmini0172-quantityreadout}
  Time readout in seconds
\end{definition}
\begin{proof}
  This is the declared model definition of time In Seconds.
\end{proof}
\begin{definition}[frequency In Hertz]
  \label{def:physics:phyx-mini-0172:phyxminiproblems-problemphyxmini0172-frequencyinhertz}
  \lean{PhyXMiniProblems.ProblemPhyXMini0172.frequencyInHertz}
  \uses{def:physics:phyx-mini-0172:phyxminiproblems-problemphyxmini0172-frequencyquantity, def:physics:phyx-mini-0172:phyxminiproblems-problemphyxmini0172-quantityreadout}
  Frequency readout in SI hertz
\end{definition}
\begin{proof}
  This is the declared model definition of frequency In Hertz.
\end{proof}
\begin{definition}[Wire Label]
  \label{def:physics:phyx-mini-0172:phyxminiproblems-problemphyxmini0172-wirelabel}
  \lean{PhyXMiniProblems.ProblemPhyXMini0172.WireLabel}
  The two supporting wires labelled in the supplied figure
\end{definition}
\begin{proof}
  This is the declared model definition of Wire Label.
\end{proof}
\begin{definition}[Load Label]
  \label{def:physics:phyx-mini-0172:phyxminiproblems-problemphyxmini0172-loadlabel}
  \lean{PhyXMiniProblems.ProblemPhyXMini0172.LoadLabel}
  The two downward loads on the support system
\end{definition}
\begin{proof}
  This is the declared model definition of Load Label.
\end{proof}
\begin{definition}[Bar Endpoint]
  \label{def:physics:phyx-mini-0172:phyxminiproblems-problemphyxmini0172-barendpoint}
  \lean{PhyXMiniProblems.ProblemPhyXMini0172.BarEndpoint}
  The endpoint of the bar at which a supporting wire is attached
\end{definition}
\begin{proof}
  This is the declared model definition of Bar Endpoint.
\end{proof}
\begin{definition}[Orientation]
  \label{def:physics:phyx-mini-0172:phyxminiproblems-problemphyxmini0172-orientation}
  \lean{PhyXMiniProblems.ProblemPhyXMini0172.Orientation}
  Qualitative orientations visible in the primary figure
\end{definition}
\begin{proof}
  This is the declared model definition of Orientation.
\end{proof}
\begin{definition}[Bar Mass Distribution]
  \label{def:physics:phyx-mini-0172:phyxminiproblems-problemphyxmini0172-barmassdistribution}
  \lean{PhyXMiniProblems.ProblemPhyXMini0172.BarMassDistribution}
  The stated mass distribution of the bar
\end{definition}
\begin{proof}
  This is the declared model definition of Bar Mass Distribution.
\end{proof}
\begin{definition}[Vibration Mode]
  \label{def:physics:phyx-mini-0172:phyxminiproblems-problemphyxmini0172-vibrationmode}
  \lean{PhyXMiniProblems.ProblemPhyXMini0172.VibrationMode}
  The vibration mode specified in the question
\end{definition}
\begin{proof}
  This is the declared model definition of Vibration Mode.
\end{proof}
\begin{definition}[Figure Feature]
  \label{def:physics:phyx-mini-0172:phyxminiproblems-problemphyxmini0172-figurefeature}
  \lean{PhyXMiniProblems.ProblemPhyXMini0172.FigureFeature}
  Individually visible, labelled features of the supplied diagram
\end{definition}
\begin{proof}
  This is the declared model definition of Figure Feature.
\end{proof}
\begin{definition}[Suspension Figure]
  \label{def:physics:phyx-mini-0172:phyxminiproblems-problemphyxmini0172-suspensionfigure}
  \lean{PhyXMiniProblems.ProblemPhyXMini0172.SuspensionFigure}
  \uses{def:physics:phyx-mini-0172:phyxminiproblems-problemphyxmini0172-wirelabel, def:physics:phyx-mini-0172:phyxminiproblems-problemphyxmini0172-barendpoint, def:physics:phyx-mini-0172:phyxminiproblems-problemphyxmini0172-figurefeature}
  Categorical evidence stored from the primary image
\end{definition}
\begin{proof}
  This is the declared model definition of Suspension Figure.
\end{proof}
\begin{definition}[Suspended Bar Wire Setup]
  \label{def:physics:phyx-mini-0172:phyxminiproblems-problemphyxmini0172-suspendedbarwiresetup}
  \lean{PhyXMiniProblems.ProblemPhyXMini0172.SuspendedBarWireSetup}
  \uses{def:physics:phyx-mini-0172:phyxminiproblems-problemphyxmini0172-lengthquantity, def:physics:phyx-mini-0172:phyxminiproblems-problemphyxmini0172-massquantity, def:physics:phyx-mini-0172:phyxminiproblems-problemphyxmini0172-forcequantity, def:physics:phyx-mini-0172:phyxminiproblems-problemphyxmini0172-timequantity, def:physics:phyx-mini-0172:phyxminiproblems-problemphyxmini0172-frequencyquantity, def:physics:phyx-mini-0172:phyxminiproblems-problemphyxmini0172-wirelabel, def:physics:phyx-mini-0172:phyxminiproblems-problemphyxmini0172-loadlabel, def:physics:phyx-mini-0172:phyxminiproblems-problemphyxmini0172-orientation, def:physics:phyx-mini-0172:phyxminiproblems-problemphyxmini0172-barmassdistribution, def:physics:phyx-mini-0172:phyxminiproblems-problemphyxmini0172-vibrationmode, def:physics:phyx-mini-0172:phyxminiproblems-problemphyxmini0172-suspensionfigure}
  All physical quantities in the experiment. Positions are fractions of the bar span measured rightward from support A: hence A has coordinate 0 and B coordinate 1. Neither beatFrequency nor either wire tension is assigned a numerical answer in this structure
\end{definition}
\begin{proof}
  This is the declared model definition of Suspended Bar Wire Setup.
\end{proof}
\begin{definition}[Matches Primary Figure]
  \label{def:physics:phyx-mini-0172:phyxminiproblems-problemphyxmini0172-matchesprimaryfigure}
  \lean{PhyXMiniProblems.ProblemPhyXMini0172.MatchesPrimaryFigure}
  \uses{def:physics:phyx-mini-0172:phyxminiproblems-problemphyxmini0172-suspendedbarwiresetup}
  The primary image shows both labelled vertical wires descending from a ceiling to opposite ends of the horizontal bar, with the cube resting on the bar. This predicate contains no force, frequency, or answer-choice value
\end{definition}
\begin{proof}
  This is the declared model definition of Matches Primary Figure.
\end{proof}
\begin{definition}[Matches Bar And Support Geometry]
  \label{def:physics:phyx-mini-0172:phyxminiproblems-problemphyxmini0172-matchesbarandsupportgeometry}
  \lean{PhyXMiniProblems.ProblemPhyXMini0172.MatchesBarAndSupportGeometry}
  \uses{def:physics:phyx-mini-0172:phyxminiproblems-problemphyxmini0172-suspendedbarwiresetup}
  Qualitative geometry and normalized load locations. Uniformity puts the bar's weight at its midpoint; the lead cube is at the stated three-quarter point from A to B
\end{definition}
\begin{proof}
  This is the declared model definition of Matches Bar And Support Geometry.
\end{proof}
\begin{definition}[Matches Stated Measurements]
  \label{def:physics:phyx-mini-0172:phyxminiproblems-problemphyxmini0172-matchesstatedmeasurements}
  \lean{PhyXMiniProblems.ProblemPhyXMini0172.MatchesStatedMeasurements}
  \uses{def:physics:phyx-mini-0172:phyxminiproblems-problemphyxmini0172-lengthincentimeters, def:physics:phyx-mini-0172:phyxminiproblems-problemphyxmini0172-massingrams, def:physics:phyx-mini-0172:phyxminiproblems-problemphyxmini0172-forceinnewtons, def:physics:phyx-mini-0172:phyxminiproblems-problemphyxmini0172-wirelabel, def:physics:phyx-mini-0172:phyxminiproblems-problemphyxmini0172-suspendedbarwiresetup}
  The numerical measurements printed in the problem: bar weight 165 N, cube weight 185 N, and, for each identical wire, length 75.0 cm and mass 5.50 g. No tension or frequency appears here
\end{definition}
\begin{proof}
  This is the declared model definition of Matches Stated Measurements.
\end{proof}
\begin{definition}[Matches Simultaneous Fundamental Pluck]
  \label{def:physics:phyx-mini-0172:phyxminiproblems-problemphyxmini0172-matchessimultaneousfundamentalpluck}
  \lean{PhyXMiniProblems.ProblemPhyXMini0172.MatchesSimultaneousFundamentalPluck}
  \uses{def:physics:phyx-mini-0172:phyxminiproblems-problemphyxmini0172-suspendedbarwiresetup}
  Both wires are plucked at their centers at the same time and are considered in their fundamental modes. This records the excitation requested by the question, without asserting the resulting beat frequency
\end{definition}
\begin{proof}
  This is the declared model definition of Matches Simultaneous Fundamental Pluck.
\end{proof}
\begin{definition}[Has Physical Parameters]
  \label{def:physics:phyx-mini-0172:phyxminiproblems-problemphyxmini0172-hasphysicalparameters}
  \lean{PhyXMiniProblems.ProblemPhyXMini0172.HasPhysicalParameters}
  \uses{def:physics:phyx-mini-0172:phyxminiproblems-problemphyxmini0172-lengthinmeters, def:physics:phyx-mini-0172:phyxminiproblems-problemphyxmini0172-massinkilograms, def:physics:phyx-mini-0172:phyxminiproblems-problemphyxmini0172-forceinnewtons, def:physics:phyx-mini-0172:phyxminiproblems-problemphyxmini0172-frequencyinhertz, def:physics:phyx-mini-0172:phyxminiproblems-problemphyxmini0172-wirelabel, def:physics:phyx-mini-0172:phyxminiproblems-problemphyxmini0172-loadlabel, def:physics:phyx-mini-0172:phyxminiproblems-problemphyxmini0172-suspendedbarwiresetup}
  Positivity and in-span conditions for the mechanical model. These are physical admissibility assumptions and do not identify an answer choice
\end{definition}
\begin{proof}
  This is the declared model definition of Has Physical Parameters.
\end{proof}
\begin{definition}[Satisfies Static Equilibrium]
  \label{def:physics:phyx-mini-0172:phyxminiproblems-problemphyxmini0172-satisfiesstaticequilibrium}
  \lean{PhyXMiniProblems.ProblemPhyXMini0172.SatisfiesStaticEquilibrium}
  \uses{def:physics:phyx-mini-0172:phyxminiproblems-problemphyxmini0172-lengthinmeters, def:physics:phyx-mini-0172:phyxminiproblems-problemphyxmini0172-forceinnewtons, def:physics:phyx-mini-0172:phyxminiproblems-problemphyxmini0172-suspendedbarwiresetup}
  Vertical force balance and moment balance about the left end of the bar. The second equality uses actual metre lever arms: the bar span multiplied by each normalized position. It is a general static-equilibrium law, not the numerical tension sought as an intermediate calculation
\end{definition}
\begin{proof}
  This is the declared model definition of Satisfies Static Equilibrium.
\end{proof}
\begin{definition}[stretched String Fundamental Hertz]
  \label{def:physics:phyx-mini-0172:phyxminiproblems-problemphyxmini0172-stretchedstringfundamentalhertz}
  \lean{PhyXMiniProblems.ProblemPhyXMini0172.stretchedStringFundamentalHertz}
  The fundamental frequency in hertz of a stretched uniform wire with SI length lengthMeters, mass massKilograms, and tension tensionNewtons: f = (1 / (2 L)) * sqrt(T / (m / L))
\end{definition}
\begin{proof}
  This is the declared model definition of stretched String Fundamental Hertz.
\end{proof}
\begin{definition}[Satisfies Stretched String Fundamental Law]
  \label{def:physics:phyx-mini-0172:phyxminiproblems-problemphyxmini0172-satisfiesstretchedstringfundamentallaw}
  \lean{PhyXMiniProblems.ProblemPhyXMini0172.SatisfiesStretchedStringFundamentalLaw}
  \uses{def:physics:phyx-mini-0172:phyxminiproblems-problemphyxmini0172-lengthinmeters, def:physics:phyx-mini-0172:phyxminiproblems-problemphyxmini0172-massinkilograms, def:physics:phyx-mini-0172:phyxminiproblems-problemphyxmini0172-forceinnewtons, def:physics:phyx-mini-0172:phyxminiproblems-problemphyxmini0172-frequencyinhertz, def:physics:phyx-mini-0172:phyxminiproblems-problemphyxmini0172-wirelabel, def:physics:phyx-mini-0172:phyxminiproblems-problemphyxmini0172-suspendedbarwiresetup, def:physics:phyx-mini-0172:phyxminiproblems-problemphyxmini0172-stretchedstringfundamentalhertz}
  Each supporting wire obeys the fundamental stretched-string law with its own static tension and its measured length and mass. The law is uniform over both labels and contains no beat-frequency answer
\end{definition}
\begin{proof}
  This is the declared model definition of Satisfies Stretched String Fundamental Law.
\end{proof}
\begin{definition}[Satisfies Beat Frequency Law]
  \label{def:physics:phyx-mini-0172:phyxminiproblems-problemphyxmini0172-satisfiesbeatfrequencylaw}
  \lean{PhyXMiniProblems.ProblemPhyXMini0172.SatisfiesBeatFrequencyLaw}
  \uses{def:physics:phyx-mini-0172:phyxminiproblems-problemphyxmini0172-frequencyinhertz, def:physics:phyx-mini-0172:phyxminiproblems-problemphyxmini0172-suspendedbarwiresetup}
  For two pure tones, the beat frequency is the absolute difference of their frequencies. This governing acoustics law relates the unknown observable to the two component frequencies but supplies no numerical result
\end{definition}
\begin{proof}
  This is the declared model definition of Satisfies Beat Frequency Law.
\end{proof}
\begin{definition}[Answer Choice]
  \label{def:physics:phyx-mini-0172:phyxminiproblems-problemphyxmini0172-answerchoice}
  \lean{PhyXMiniProblems.ProblemPhyXMini0172.AnswerChoice}
  Labels of the four frequency choices printed with the problem
\end{definition}
\begin{proof}
  This is the declared model definition of Answer Choice.
\end{proof}
\begin{definition}[frequency Hertz]
  \label{def:physics:phyx-mini-0172:phyxminiproblems-problemphyxmini0172-answerchoice-frequencyhertz}
  \lean{PhyXMiniProblems.ProblemPhyXMini0172.AnswerChoice.frequencyHertz}
  \uses{def:physics:phyx-mini-0172:phyxminiproblems-problemphyxmini0172-answerchoice}
  The hertz value displayed beside each answer label
\end{definition}
\begin{proof}
  This is the declared model definition of frequency Hertz.
\end{proof}
\begin{definition}[recorded Answer Choice]
  \label{def:physics:phyx-mini-0172:phyxminiproblems-problemphyxmini0172-recordedanswerchoice}
  \lean{PhyXMiniProblems.ProblemPhyXMini0172.recordedAnswerChoice}
  \uses{def:physics:phyx-mini-0172:phyxminiproblems-problemphyxmini0172-answerchoice}
  Dataset metadata: the recorded answer is A; this is never a premise
\end{definition}
\begin{proof}
  This is the declared model definition of recorded Answer Choice.
\end{proof}
\begin{definition}[Is Closest Answer Choice]
  \label{def:physics:phyx-mini-0172:phyxminiproblems-problemphyxmini0172-isclosestanswerchoice}
  \lean{PhyXMiniProblems.ProblemPhyXMini0172.IsClosestAnswerChoice}
  \uses{def:physics:phyx-mini-0172:phyxminiproblems-problemphyxmini0172-frequencyquantity, def:physics:phyx-mini-0172:phyxminiproblems-problemphyxmini0172-frequencyinhertz, def:physics:phyx-mini-0172:phyxminiproblems-problemphyxmini0172-answerchoice}
  A displayed answer is uniquely closest to the model's beat frequency
\end{definition}
\begin{proof}
  This is the declared model definition of Is Closest Answer Choice.
\end{proof}
\begin{lemma}[stated Wire Measurements In SI]
  \label{lem:physics:phyx-mini-0172:phyxminiproblems-problemphyxmini0172-statedwiremeasurementsinsi}
  \lean{PhyXMiniProblems.ProblemPhyXMini0172.statedWireMeasurementsInSI}
  \uses{def:physics:phyx-mini-0172:phyxminiproblems-problemphyxmini0172-lengthinmeters, def:physics:phyx-mini-0172:phyxminiproblems-problemphyxmini0172-massinkilograms, def:physics:phyx-mini-0172:phyxminiproblems-problemphyxmini0172-wirelabel, def:physics:phyx-mini-0172:phyxminiproblems-problemphyxmini0172-suspendedbarwiresetup, def:physics:phyx-mini-0172:phyxminiproblems-problemphyxmini0172-matchesstatedmeasurements}
  The original 75.0 cm and 5.50 g data have SI readouts 3/4 m and 11/2000 kg. This is a unit-conversion intermediate, not a frequency conclusion
\end{lemma}
\begin{proof}
  The conclusion follows from the governing relations and figure readouts named in the statement of stated Wire Measurements In SI.
\end{proof}
\begin{lemma}[support Tensions In Newtons]
  \label{lem:physics:phyx-mini-0172:phyxminiproblems-problemphyxmini0172-supporttensionsinnewtons}
  \lean{PhyXMiniProblems.ProblemPhyXMini0172.supportTensionsInNewtons}
  \uses{def:physics:phyx-mini-0172:phyxminiproblems-problemphyxmini0172-forceinnewtons, def:physics:phyx-mini-0172:phyxminiproblems-problemphyxmini0172-suspendedbarwiresetup, def:physics:phyx-mini-0172:phyxminiproblems-problemphyxmini0172-matchesbarandsupportgeometry, def:physics:phyx-mini-0172:phyxminiproblems-problemphyxmini0172-matchesstatedmeasurements, def:physics:phyx-mini-0172:phyxminiproblems-problemphyxmini0172-hasphysicalparameters, def:physics:phyx-mini-0172:phyxminiproblems-problemphyxmini0172-satisfiesstaticequilibrium}
  Static equilibrium gives 128.75 N in wire A and 221.25 N in wire B. These are derived intermediate tensions; neither is the requested beat frequency
\end{lemma}
\begin{proof}
  The conclusion follows from the governing relations and figure readouts named in the statement of support Tensions In Newtons.
\end{proof}
% --- Archon declaration topology end ---
% --- Archon physics formalization source end ---
